% latex-header.tex — md-convert 的 XeLaTeX 导言（便携版）
%
% 用法（手动调用 pandoc 时）：
%   pandoc input.md -o output.pdf --pdf-engine=xelatex -H latex-header.tex \
%       -V geometry:top=2.54cm -V geometry:bottom=2.54cm \
%       -V geometry:left=3.17cm -V geometry:right=3.17cm -V linestretch=1.5
%
% 说明：scripts/to_pdf.py 会动态生成等价导言，
% 并在检测到本机字体文件时用 Path= 方式精确加载 Heavy 字体；
% 本文件供手动调用或自定义时使用，仅按“字体族名”查找已安装字体。

% ---- 西文 / 数字 -----------------------------------------------------------
\setmainfont{Times New Roman}

% ---- 中文 -------------------------------------------------------------------
\usepackage{xeCJK}
\setCJKmainfont{SimSun}[AutoFakeBold=2.5]

% ---- 强调字重：思源宋体 Heavy（无伪粗体） ------------------------------------
% 若本机未安装该字体族，回退为宋体 AutoFakeBold。
\IfFontExistsTF{Noto Serif SC Heavy X}{%
  \setCJKfamilyfont{emphsong}{Noto Serif SC Heavy X}%
}{%
  \setCJKfamilyfont{emphsong}{SimSun}[AutoFakeBold=2.5]%
}
\newcommand{\key}[1]{{\CJKfamily{emphsong}\bfseries #1}}

% ---- 数学字体：与正文（Times 系）同风格 --------------------------------------
% pandoc 的 xelatex 模板已加载 unicode-math；按可用性回退。
\IfFontExistsTF{STIX Two Math}{%
  \setmathfont{STIX Two Math}%
}{%
  \IfFontExistsTF{TeX Gyre Termes Math}{%
    \setmathfont{TeX Gyre Termes Math}%
  }{}%
}

% 本机有字体文件但未安装到系统时，把上面的 \IfFontExistsTF 块替换为：
%   \setCJKfamilyfont{emphsong}{NotoSerifSC-HeavyX.ttf}%
%     [Path=C:/path/to/fonts/, BoldFont=NotoSerifSC-HeavyX.ttf]
